%%%%%%%%%%%%%%%%%%%%%%%%%%%%%%%%%%%%%%%%%%%%%%%%%%%%%%%%%%%%%%%%%%%%%%
% Lớp 10 | Xác suất | Xác suất của biến cố
% Dạng 0D0?2-7: Liên quan việc đếm số
% Số câu: 24  (sinh tự động bởi tach_cau.py; từ nay thầy sửa trực tiếp file này)
%%%%%%%%%%%%%%%%%%%%%%%%%%%%%%%%%%%%%%%%%%%%%%%%%%%%%%%%%%%%%%%%%%%%%%

%%%%%%%%%%%%%%%% TRẮC NGHIỆM 4 PHƯƠNG ÁN %%%%%%%%%%%%%%%%

% nguồn: MaHoaNH_TracNghiem_Lop10
\begin{ex}%[0D0H2-7]%[ma:4f58bc31]
	Từ các chữ số $1$; $2$; $4$; $6$; $8$; $9$ lấy ngẫu nhiên một số. Xác suất  để lấy được một số nguyên tố là?
	\choice 
		{$\dfrac{1}{2}$} 
		 {$\dfrac{1}{3}$}
		{$\dfrac{1}{4}$} 
		{\True $\dfrac{1}{6}$} 
	 \loigiai{
		Gọi $A\colon$ \lq\lq  Lấy được một số nguyên tố\rq\rq.\\
		Khi đó $A = \left\lbrace 2\right\rbrace \Rightarrow n(A) = 1$ và $n(\Omega) = 6$. Do  đó, xác suất để lấy được một số nguyên tố là
		$$ \mathrm{P}(A) = \dfrac{n(A)}{n(\Omega)} = \dfrac{1}{6}.  $$
	} 
\end{ex}

% nguồn: MaHoaNH_TracNghiem_Lop10
\begin{ex}%[0D0V2-7]%[ma:4eaef662]
	Gọi $S$ là tập hợp số có $5$ chữ số khác nhau được lập  từ các số $0$, $1$, $2$, $3$, $4$, $5$, $6$, $7$. Chọn $1$ số từ $S$, tính xác suất sao cho số $5$ và $1$ không đứng cạnh nhau?
	 \choice
		{$\dfrac{43}{49}$} 
		 {$\dfrac{6}{49}$}
		 {\True $\dfrac{44}{49}$}
		 {$\dfrac{5}{49}$}
	\loigiai{ 
		Số  phần tử của không gian mẫu là $n(\Omega)=\mathrm{A}_{8}^{5}-\mathrm{A}_{7}^{4}=5\,880$.\\
		Gọi $A$ là  biến cố \lq\lq tập hợp số có $5$ chữ số khác nhau mà số $5$ và $1$ không đứng cạnh nhau\rq\rq.\\
		Khi đó $\overline{A}$ là biến cố \lq\lq tập hợp số có $5$ chữ  số khác nhau mà số $5$ và $1$ đứng cạnh nhau\rq\rq.\\
		Ta có các  trường hợp thuận lợi cho biến cố $\overline{A}$ là
		\begin{center}
			\begin{tabular}{|c|c|c|c|c|} \hline
				$1$& $2$& $3$& $4$&  $5$ \\ \hline
			\end{tabular}
		\end{center}
		Trường hợp $1$: số $5$ và $1$ xếp vào vị trí số thứ nhất và thứ hai, $3$ số còn lại chọn vào $3$ vị trí còn lại có số cách chọn là $\mathrm{A}_{6}^{3}$. Trường hợp này  có số cách lập số là $2\cdot \mathrm{A}_{6}^{3}$.\\
		Trường hợp $2$: số $5$ và $1$ xếp vào vị trí $\left\{2,3\right\}$, $\left\{3,4\right\}$, $\left\{4, 5\right\}$. Khi đó $3$ số còn lại  chọn vào $3$ vị trí còn, nhưng loại đi những chữ số $0$ ở vị trí $1$;.\\
		Mỗi cặp số $5$ và $1$ xếp vào vị trí $\left\{2;3\right\}$, số  cách chọn là $2\cdot (\mathrm{A}_{6}^{3}-\mathrm{A}_{5}^{2})=200$.\\
		Tương tự các cặp $\left\{3;4\right\}$, $\left\{4;5\right\}$ có số cách $2\cdot 2\cdot  (\mathrm{A}_{6}^{3}-\mathrm{A}_{5}^{2})=400$.\\
		Số cách chọn của  $\overline{A}$ là $n(\overline{A})=200+400=600$.\\
		Xác suất  cần tìm $\mathrm{P}(A)=1-\mathrm{P}(\overline{A})=1-\dfrac{600}{5\,880}=\dfrac{44}{49}$.
	 }
\end{ex}

% nguồn: MaHoaNH_TracNghiem_Lop10
\begin{ex}%[0D0V2-7]%[ma:91740cb7]
	Gọi $S$ là tập hợp các số tự nhiên gồm bốn chữ số phân biệt được chọn từ  các chữ số $0$, $1$, $2$, $3$, $4$, $5$, $6$, $7$. Chọn ngẫu nhiên một số từ $S$. Xác suất để chọn được số không vượt quá $2\,020$ là
	 \choice
		{\True  $\dfrac{43}{294}$}
		 {$\dfrac{239}{294}$}
		{$\dfrac{6}{7}$} 
		{$\dfrac{36}{49}$} 
	 \loigiai{
		Số phần tử của  tập hợp $S$ là $n(S)=7\cdot \mathrm{A}_7^3=1\,470$.\\
		Số phần tử của không gian mẫu là  $n\left(\Omega\right)=\mathrm{C}_{1\,470}^1=1\,470$.\\
		Gọi  biến cố $A$: \lq\lq số được chọn không vượt quá $2\,020$\rq\rq.\\
		Suy ra biến cố  đối $\overline{A}$: \lq\lq số được chọn lớn hơn $2\,020$\rq\rq.\\
		Giả sử $n=\overline{abcd}\in  \overline{A}$ và $n>2\,020$.
		\begin{description}
			\item[Trường hợp 1:] $a=2$; $b=0$  thì $c\in \{3;4;5;6;7\}$ nên $c$ có $5$ cách chọn và $d$ có $5$ cách chọn.\\
			Suy ra có $1\cdot  1\cdot 5\cdot 5=25$ số.
			\item[Trường hợp 2:] $a=2$; $b\in \{1;3;4;5;6;7\}$ thì $\overline{cd}$ có $\mathrm{A}_6^2$ cách chọn và sắp xếp.\\ 
			 Suy ra có $1\cdot 6\cdot \mathrm{A}_6^2=180$ số.
			\item[Trường hợp 3:] $a\in \{3;4;5;6;7\}$ thì  $\overline{bcd}$ có $\mathrm{A}_7^3$ cách chọn và sắp xếp.\\
			 Suy ra có $5\cdot \mathrm{A}_7^3=1\,050$ số.
		\end{description}
		 Số phần tử của biến cố $\overline{A}$ là $n\left(\overline{A}\right)=25+180+1\,050=1\,255$.\\
		Suy  ra $\mathrm{P}\left(\overline{A}\right)=\dfrac{1\,255}{1\,470}=\dfrac{251}{294}$.\\
		Vậy $\mathrm{P}(A)=1-\mathrm{P}\left(\overline{A}\right)=1-\dfrac{251}{294}=\dfrac{43}{294}$. 
	 }
\end{ex}

% nguồn: MaHoaNH_TracNghiem_Lop10
\begin{ex}%[0D0V2-7]%[ma:f16c7406]
	Gọi $S$ là tập hợp tất cả các số có  $5$ chữ số khác nhau được lập từ các chữ số $0, 1, 2, 3, 4, 5, 6, 7$. Chọn ngẫu nhiên một số từ $S$. Tính xác suất để số chọn được chia hết cho $5$, luôn có mặt các chữ số $2, 3, 4$ và chúng đứng cạnh nhau.
	\choice 
		 {$\dfrac{1}{140}$}
		 {$\dfrac{1}{392}$}
		 {$\dfrac{4}{245}$}
		 {\True $\dfrac{3}{196}$}
	\loigiai{ 
		Gọi số cần tìm là $\overline{abcde}$. Trong đó  các chữ số $a, b, c, d, e$ đôi một khác nhau và thuộc tập hợp $M=\{0, 1, 2, 3, 4, 5, 6, 7 \}$, $a\ne 0$.\\
		Số các chữ số có $5$ chữ  số đôi một khác nhau thuộc tập hợp $M$ là $n\left(\Omega\right)= 7\cdot\mathrm{A}_7^4=5\,880$.\\
		Gọi $A$ là biến cố \lq\lq  số chọn được chia hết cho $5$, luôn có mặt các chữ số $2, 3, 4$ và chúng đứng cạnh nhau\rq\rq.\\
		Vì $\overline{abcde}\vdots 5$ nên  $e\in\{0, 5\}$.\\
		Vì các chữ số $2, 3, 4$ luôn có mặt  và chúng đứng cạnh nhau nên xếp $3$ ba chữ số $2, 3, 4$ thành một nhóm $(234)$ có $3!$ cách xếp.
		\begin{itemize}
			\item {\bf Trường hợp 1:} $e=0$.\\ 
			Xếp chữ số  $e$ có $1$ cách.\\
			Xếp nhóm các chữ số $(234)$  có $2\cdot3!=12$ cách.\\
			Xếp các chữ số  $1, 5, 6, 7$ vào vị trí còn lại có $4$ cách.\\
			Vậy  có $12\cdot4=48$ cách.
			\item  {\bf Trường hợp 2:} $e=5$.
			\begin{itemize}
				\item Xếp nhóm các chữ số $(234)$ vào vị trí  $(abc)$ có $3!=6$ cách.\\
				Xếp  các chữ số $0, 1, 6, 7$ vào vị trí $d$ có $4$ các xếp.\\
				Trường hợp này có $6\cdot4=24$  cách.
				\item Xếp nhóm các  chữ số $(234)$ vào vị trí $(bcd)$ có $3!=6$ cách.\\
				Xếp các chữ số $1, 6, 7$ vào vị trí $a$ có $3$ các  xếp.\\
				Trường hợp này có  $6\cdot3=18$ cách.
			\end{itemize}
			Vậy có $24 +  18 = 42$ cách.
		\end{itemize}
		Do đó số kết quả thuận lợi của biến cố $A$ là  $n\left(A\right)= 48+42 = 90$.\\
		Xác suất của biến cố  $A$ là
		 $\mathrm{P}(A)=\dfrac{n(A)}{n\left(\Omega\right)}=\dfrac{90}{5\,880}=\dfrac{3}{196}$.
	 }
\end{ex}


%%%%%%%%%%%%%%%% ĐÚNG/SAI %%%%%%%%%%%%%%%%

% nguồn: MaHoaNH_DungSai_Lop10
\begin{ex}%[0D0H2-7]%[ma:b11f587a]
	Một chiếc hộp chứa $20$ chiếc thẻ được đánh số liên tiếp từ $1$ đến $20$. Lấy ngẫu nhiên hai thẻ trong  hộp. Khi đó
	 \choiceTF
		{\True  Số phần tử của không gian mẫu $n(\Omega)=190$}
		{\True Số cách để bốc được hai thẻ mang số  chẵn là $n(B)=\mathrm{C}_{10}^2=45$ cách}
		{\True Xác suất để hai thẻ lấy được một thẻ mang số chẵn, một thẻ  mang số lẻ là $\mathrm{P}(C)=\dfrac{10}{19}$}
		{ Số cách để lấy hai thẻ đều là là số nguyên tố lẻ  là $n(D)=\mathrm{C}_{8}^2$}
	\loigiai{ 
		\begin{itemchoice}
			\itemch  Vì số phần tử của không gian mẫu $n(\Omega)=\mathrm{C}_{20}^2=190$.
			\itemch  Số cách để bốc được hai thẻ cùng là số chẵn là $n(B)=\mathrm{C}_{10}^2=45$.
			\itemch Xác suất để lấy một thẻ là số chẵn, một thẻ  là số lẻ là $\mathrm{P}(C)=\dfrac{10}{19}$.\\
			Có $10$ cách lấy một thẻ là số chẵn, có  $10$ cách lấy một thẻ là số lẻ.\\
			Vậy $n(C)=10.10=100$ nên  $\mathrm{P}(C)=\dfrac{n(C)}{n(\Omega)}=\dfrac{100}{190}=\dfrac{10}{19}$.
			\itemch Số cách để lấy hai thẻ đều là là số nguyên tố lẻ là  $n(D)=\mathrm{C}_{8}^2$.\\
			Các số nguyên tố lẻ là $3; 5; 7; 11;  13; 17; 19$. \\Vậy $n(D)=\mathrm{C}_7^2$.
	\end{itemchoice}}
\end{ex}

% nguồn: MaHoaNH_DungSai_Lop10
\begin{ex}%[0D0H2-7]%[ma:a1ae49c7]
	Một chiếc hộp gồm có  $9$ thẻ được đánh số liên tiếp từ $1$ đến $9$. Lấy ngẫu nhiên hai thẻ trong hộp. Khi đó
	\choiceTF 
		{\True Số phần tử  của không gian mẫu $n(\Omega)=36$}
		{ Số  cách để bốc được hai thẻ mang số chẵn là $n(B)=10$ cách}
		{\True Xác suất để hai thẻ lấy được có tích của chúng là số lẻ là $\mathrm{P}(C)=\dfrac{5}{18}$} 
		{ Xác suất để hai thẻ lấy được có tích của chúng  là số chẵn là $\mathrm{P}(D)=\dfrac{1}{6}$}
	\loigiai{ 
		\begin{itemchoice}
			\itemch Vì chọn $2$  số trong $9$ số là $n(\Omega)=\mathrm{C}_{9}^2=36$.
			\itemch Vì chọn hai số chẵn thuộc  tập $\{2; 4; 6; 8\}$ có $n(B)=\mathrm{C}_{4}^2=6$ cách.
			\itemch Vì chọn  hai số lẻ thuộc tập $\{1; 3; 5; 7; 9\}$ có $n(C)=\mathrm{C}_5^2=10$ cách.\\
			Xác suất để chọn được hai số  lẻ là $\mathrm{P}(C)=\dfrac{n(C)}{n(\Omega)}=\dfrac{10}{36}=\dfrac{5}{18}$.
			\itemch Gọi $D$ là biến cố: \lq\lq Tích của  số ghi trên hai thẻ là một số chẵn\rq\rq .\\
			 Tập các số chẵn là $\{2;4;6;8\}$; tập các số lẻ là $\{1;3;5;7;9\}$.
			\begin{itemize}
				\item \textbf{Trường hợp  1.} Hai số chẵn: $\mathrm{C}_{4}^2=6$.
				\item \textbf{Trường hợp 2.}  Một lẻ, một chẵn: $\mathrm{C}_5^1 \cdot \mathrm{C}_{4}^1=20$
			\end{itemize}
			Suy ra $ n(D)=6+20=26 \Rightarrow \mathrm{P}(D)=\dfrac{n(D)}{n(\Omega)}=\dfrac{26}{36}=\dfrac{13}{18}$. 
	\end{itemchoice}}
\end{ex}

% nguồn: MaHoaNH_DungSai_Lop10
\begin{ex}%[0D0H2-7]%[ma:218c41ba]
	Gọi $A$ là tập hợp các số tự nhiên có $2$ chữ số nhỏ hơn $100$. Lấy ra $1$ số tự nhiên bất kỳ  trong $A$. Khi đó
	 \choiceTF
		{\True Không gian mẫu $\Omega=\{10;11;12;13; \ldots; 98;99\}$ nên ta có  $n(\Omega)=90$}
		{\True Gọi $A$ là biến cố \lq\lq Lấy được  số tự nhiên lẻ\rq\rq\ nên ta có $\mathrm{P}(A)=0{,}5$}
		{Gọi $B$ là biến cố \lq\lq Lấy được số tự nhiên chia hết  cho $3$\rq\rq\ nên ta có $\mathrm{P}(B)=\dfrac{4}{9}$}
		{Gọi $C$ là biến cố \lq\lq Lấy được  số có hai chữ số giống nhau\rq\rq\ nên ta có $\mathrm{P}(C)=\dfrac{1}{10}$}
	\loigiai{ 
		\begin{itemchoice}
			\itemch Không gian mẫu $\Omega=\{10;11;12;13; \ldots; 98;99\}$.\\ 
			Số phần tử  không gian mẫu là $n(\Omega)=99-10+1=90$.
			\itemch Gọi $A$ là biến cố \lq\lq Lấy được số  tự nhiên lẻ\rq\rq.\\
			Khi đó  $A=\{11;13;15;17;19; \ldots; 99\}$, suy ra $n(A)=[(99-11): 2]+1=45$.\\
			 Vậy $\mathrm{P}(A)=\dfrac{45}{90}=0{,}5$.
			\itemch Gọi $B$ là biến cố \lq\lq Lấy  được số tự nhiên chia hết cho $3$\rq\rq.\\
			Khi  đó $B=\{12;15;18; \ldots; 99\}$, suy ra $n(B)=[(99-12): 3]+1=30$.\\
			Vậy  $\mathrm{P}(B)=\dfrac{30}{90}=\dfrac{1}{3}$.
			\itemch Gọi $C$ là biến cố \lq\lq Lấy được số có hai chữ  số giống nhau\rq\rq.\\
			Khi đó $C=\{11;22;33;44; \ldots;  99\}$, suy ra $n(C)=9$.\\
			Vậy  $\mathrm{P}(C)=\dfrac{9}{90}=\dfrac{1}{10}$.
		\end{itemchoice}
	 }
\end{ex}

% nguồn: MaHoaNH_DungSai_Lop10
\begin{ex}%[0D0H2-7]%[ma:8b211d80]
	Chọn ngẫu nhiên một số nguyên dương nhỏ hơn  $9$. Các khẳng định sau đúng hay sai.
	 \choiceTF
		{\True Xác suất của  biến cố chọn được số $8$ là $\dfrac{1}{8}$}
		 {\True Xác suất của biến cố chọn được số $0$ là $0$}
		{\True Xác suất của biến cố chọn được  một số nguyên tố là $\dfrac{1}{2}$}
		{Xác suất của biến cố lấy  được một hợp số là $\dfrac{4}{8}$}
	\loigiai{ 
		\begin{itemchoice}
			\itemch Số phần  tử của không gian mẫu là $n(\Omega)=8$.\\
			Gọi $ A $ là biến cố \lq\lq chọn được số $8$\rq\rq  $\Rightarrow n(A)=1$.\\
			Vậy  $\mathrm{P}(A)=\dfrac{n(A)}{n(\Omega)}=\dfrac{1}{8}$.
			\itemch  Số phần tử của không gian mẫu là $n(\Omega)=8$.\\
			Gọi $A$ là biến cố \lq\lq chọn được  số $0$\rq\rq $\Rightarrow n(A)=0$.\\
			Vậy  $\mathrm{P}(A)=\dfrac{n(A)}{n(\Omega)}=0$.
			\itemch Số phần tử  của không gian mẫu là $n(\Omega)=8$.\\
			Gọi $ A $ là biến cố \lq\lq chọn  được một số nguyên tố\rq\rq $\Rightarrow n(A)=4$.\\
			Vậy $\mathrm{P}(A)=\dfrac{n(A)}{n(\Omega)}=\dfrac{1}{2}$. 
			\itemch  Số phần tử của không gian mẫu là $n(\Omega)=8$.\\
			Gọi $ A  $ là biến cố \lq\lq chọn được một hợp số\rq\rq $\Rightarrow n(A)=3$.\\
			Vậy  $\mathrm{P}(A)=\dfrac{n(A)}{n(\Omega)}=\dfrac{3}{8}$.
		\end{itemchoice}
	 }
\end{ex}

% nguồn: MaHoaNH_DungSai_Lop10
\begin{ex}%[0D0H2-7]%[ma:f15ca3d4]
	 Trong một hòm phiếu có $9$ lá phiếu ghi các số tự nhiên từ $1$ đến $9$ (mỗi lá ghi một số, không có hai lá phiếu nào được ghi cùng một số). Rút ngẫu nhiên cùng lúc hai lá phiếu. Khi đó
	 \choiceTF
		{Số  phần tử của không gian mẫu là $72$}
		{\True Xác suất của biến cố $A$: \lq\lq Hai phiếu rút được đều là  số lẻ\rq\rq\, là $\dfrac{5}{18}$}
		{ Xác suất của biến cố $B$:  \lq\lq Tổng hai số ghi trên hai lá phiếu rút được là một số chẵn\rq\rq\, là $\dfrac{1}{6}$}
		 {\True Xác suất của biến cố $C$: \lq\lq Tổng hai số ghi trên hai lá phiếu rút được là một số lẻ lớn hơn hoặc bằng $15$\rq\rq\, là $\dfrac{1}{12}$}
	 \loigiai{
		\begin{itemchoice}
			\itemch Số phần tử của không gian mẫu là $n(\Omega)=\mathrm{C}_9^2=36$. 
			\itemch Hai phiếu rút được đều  là số lẻ, có $\mathrm{C}_5^2=10$ (cách chọn) $\Rightarrow n(A)=10$.\\
			Vậy xác suất của  biến cố $A$ là $\mathrm{P}(A)=\dfrac{n(A)}{n(\Omega)}=\dfrac{10}{36}=\dfrac{5}{18}$.
			 \itemch $B$: \lq\lq Tổng hai số ghi trên hai phiếu rút được là chẵn\rq\rq\, nên có hai khả năng xảy ra:

			\begin{itemize}
				\item Cả hai số ghi trên hai  phiếu rút được đều là số lẻ, có $\mathrm{C}_5^2=10$ (cách chọn).
				\item Cả hai số ghi trên hai phiếu rút được đều là số chẵn, có  $\mathrm{C}_{4}^2=6$ (cách chọn).
			\end{itemize}
			Suy  ra $ n(B)=10+6=16$.\\
			Vậy xác  suất của biến cố $B$ là $\mathrm{P}(B)=\dfrac{n(B)}{n(\Omega)}=\dfrac{16}{36}=\dfrac{4}{9}$.
			\itemch $A=$ \lq\lq Tổng hai số ghi trên hai lá phiếu rút được là một số  lẻ lớn hơn hoặc bằng $15$\rq\rq .\\
			Ta có các cặp số có tổng là số lẻ và lớn hơn hoặc bằng $15$ là $(6; 9)$; $(7; 8)$; $(9; 7) \Rightarrow n(A)=3.$ 
			Vậy xác suất của biến cố $A$ là  $\mathrm{P}(A)=\dfrac{3}{36}=\dfrac{1}{12}$.
	\end{itemchoice}}
\end{ex}

% nguồn: MaHoaNH_DungSai_Lop10
\begin{ex}%[0D0H2-7]%[ma:c851d204]
	Gọi $X$ là tập hợp các số  tự nhiên có $2$ chữ số khác nhau tạo nên từ các chữ số $0,2,3,8,9$. Lấy ngẫu nhiên một số từ tập $X$.
	\choiceTF 
		{\True Xác suất của biến cố  lấy được chữ số đầu tiên không nhỏ hơn $8$ (chữ số đầu tiên là chữ số hàng chục) là $\dfrac{1}{2}$}
		{Xác suất của biến cố lấy được chữ số đầu tiên nhỏ hơn $3$ (chữ số  đầu tiên là chữ số hàng chục) là $\dfrac{1}{3}$}
		 {\True Xác suất của biến cố lấy được số lẻ là $\dfrac{3}{8}$}
		{ Xác suất của biến cố lấy được số chẵn là $\dfrac{9}{16}$} 
	\loigiai{ 
		\begin{itemchoice}
			\itemch Số phần tử của không gian mẫu  là $n(\Omega)=4\cdot 4=16$.\\
			Gọi $ A $ là biến cố  \lq\lq lấy được chữ số đầu tiên không nhỏ hơn $8$\rq\rq.\\
			Gọi số thỏa yêu cầu bài toán  là $\overline{ab}$ $\Rightarrow n(A)=2\cdot 4=8$.\\
			 Vậy $\mathrm{P}(A)=\dfrac{n(A)}{n(\Omega)}=\dfrac{1}{2}$.
			\itemch Số phần tử của không gian mẫu là $n(\Omega)=4\cdot  4=16$.\\
			Gọi $  A $ là biến cố \lq\lq lấy được có chữ số đầu tiên nhỏ hơn $3$\rq\rq.\\
			Gọi số thỏa yêu  cầu bài toán là $\overline{ab}$ $\Rightarrow n(A)=1\cdot 4=4$.\\
			Vậy  $\mathrm{P}(A)=\dfrac{n(A)}{n(\Omega)}=\dfrac{1}{4}$.
			\itemch Số  phần tử của không gian mẫu là $n(\Omega)=4\cdot 4=16$.\\
			Gọi $ A $ là biến cố \lq\lq lấy được  số lẻ\rq\rq.\\
			Gọi số thỏa yêu cầu bài toán là $\overline{ab}$ $\Rightarrow n(A)=2\cdot  3=6$.\\
			 Vậy $\mathrm{P}(A)=\dfrac{n(A)}{n(\Omega)}=\dfrac{3}{8}$.
			\itemch Số phần tử của không gian mẫu là  $n(\Omega)=4\cdot 4=16$.\\
			Gọi $ A $ là biến  cố \lq\lq lấy được số chẵn\rq\rq.\\
			Gọi  số thỏa yêu cầu bài toán là $\overline{ab}$ $\Rightarrow n(A)=1\cdot 4+2\cdot 3=10$.\\
			 Vậy $\mathrm{P}(A)=\dfrac{n(A)}{n(\Omega)}=\dfrac{5}{8}$.
		\end{itemchoice}
	 }
\end{ex}


%%%%%%%%%%%%%%%% TRẢ LỜI NGẮN %%%%%%%%%%%%%%%%

% nguồn: MaHoaNH_TraLoiNgan_Lop10
\begin{bt}%[0D0H2-7]%[ma:c439b5b6]
	Chọn ngẫu nhiên $2$ số trong tập hợp $X=\{1;  2; 3; \ldots; 50\}$. Tính xác suất của biến cố sau
	$A$: \lq\lq Hai số được chọn là số  chẵn\rq\rq. (kết quả làm tròn đến hàng phần trăm)
	\shortans{0,24} 
	 \loigiai{
		Số phần tử của không gian mẫu  là $n(\Omega)=\mathrm{C}_{50}^2=1\,225$.\\
		Trong tập hợp $X$ có $25$ số chẵn $\{2; 4; 6 \ldots; 50\}$, nên số  cách lấy ra 2 số chẵn là $\mathrm{C}_{25}^2=300$ (cách). Do đó, $n(A)=300$.\\
		Xác suất của biến  cố $A$ là $\mathrm{P}(A)=\dfrac{n(A)}{n(\Omega)}=\dfrac{300}{1225}=\dfrac{12}{49}\approx0{,}24$.\\
	 }
\end{bt}

% nguồn: MaHoaNH_TraLoiNgan_Lop10
\begin{bt}%[0D0H2-7]%[ma:4237f440]
	Một hộp đựng $11$ quả cầu được đánh số từ $1$ đến $11$. Có bao nhiêu cách chọn ra  $3$ quả cầu thỏa mãn tổng ba số ghi trên $3$ quả cầu chia hết cho $3$?
	\shortans{57} 
	\loigiai{ 
		Từ $1$ đến $11$ có $3$ số chia hết  cho $3$ là $A=\{3;6;9\}$; $4$ số chia $3$ dư $1$ là tập $B=\{1;4;7;10\}$; $4$ số chia $3$ dư $2$ là tập $C=\{2;5;8;11\}$.\\
		Chọn ra $3$ quả cầu có tổng chia hết cho $3$, ta  có các trường hợp sau
		\begin{enumerate}[TH 1.]
			 \item Chọn $3$ quả cầu từ tập $A$, có số cách chọn là $1$.
			\item Chọn $3$  quả cầu từ tập $B$, có số cách chọn là $\mathrm{C}_4^3=4$.
			\item Chọn $3$ quả cầu từ tập $C$, có số  cách chọn là $\mathrm{C}_4^3=4$.
			\item Chọn $1$ quả cầu từ hộp $A$, $1$ quả cầu từ hộp $B$, $1$ quả  cầu từ hộp $C$, có số cách chọn là $3\cdot 4\cdot 4=48$.
		\end{enumerate}
		Vậy  số cách chọn thỏa mãn bài toán là $1+4+4+48=57$.
	 }
\end{bt}

% nguồn: MaHoaNH_TraLoiNgan_Lop10
\begin{bt}%[0D0H2-7]%[ma:e116425d]
	Cho tập hợp  $A=\{0;1;2;3;4;6;8\}$. Có thể lập được bao nhiêu số chẵn có $5$ chữ số đôi một khác nhau từ các chữ số thuộc $A$?
	 \shortans{1560}
	\loigiai{ 
		Xét  các số thỏa mãn bài toán, ta có hai trường hợp
		\begin{itemize}
			\item Số có dạng  $\overline{abcd0}$.\\
			 Có $\mathrm{A}_6^4$ cách chọn chữ số và xếp vào $\overline{abcd}$.\\
			Vậy trường hợp này có $\mathrm{A}_6^4$ số. 
			\item  Số có dạng $\overline{abcde}$, $e\in\{2;4;6;8\}$.\\
			Có $4$ cách chọn chữ số cho $e$.\\ 
			Có $5$ cách  chọn chữ số cho $a$.
			Có $\mathrm{A}_5^3$ cách chọn chữ số và  xếp vào $\overline{bcd}$.\\
			 Vậy trường hợp này có $4\cdot5\cdot\mathrm{A}_5^3$ số.
		\end{itemize}
	Vậy số các số  cần tìm là $S=\mathrm{A}_6^4+4\cdot5\cdot\mathrm{A}_5^3=1560$ số.}
\end{bt}

% nguồn: MaHoaNH_TraLoiNgan_Lop10
\begin{bt}%[0D0H2-7]%[ma:493bc8b7]
	Chọn ngẫu nhiên một số tự nhiên có $4$ chữ số được lập thành từ các số $0 $; $1$; $2$; $3$; $4$; $5$. Tính xác suất để số được chọn  chia hết cho $3$ (\textit{kết quả làm tròn đến chữ số thập phân thứ hai sau dấu phẩy}).
	\shortans{0,33} 
	\loigiai{ 
		Số các số $4$ chữ số được lập thành từ các  số $0$; $1$; $2$; $3$; $4$; $5$ là $5\cdot 6\cdot 6\cdot 6=1\,080$.\\
		Suy ra  $n(\Omega)=\mathrm{C}_{1\,080}^1=1\,080$.\\
		Gọi biến cố $A$: \lq\lq  Số được chọn chia hết cho $3$\rq\rq.\\
		Số các số $3$ chữ số được lập thành từ các số $0$; $1$; $2$; $3$; $4$;  $5$ là $5\cdot 6\cdot6=180$.\\
		Trong $180$ số đó  có $x$ số chia hết cho $3$; $y$ số chia $3$ dư $1$ và có $z$ số chia $3$ dư $2$.
		\begin{itemize}
			\item Với mỗi số chia hết cho $3$ thì ta có $2$ cách thêm số vào cuối để được $1$ số có $4$ chữ số và chia hết cho $3$, các số có thể thêm vào là  $0$ hoặc $3$.
			\item Với mỗi số chia $3$ dư $1$ thì ta có $2$ cách thêm số vào  cuối để được $1$ số có $4$ chữ số và chia hết cho $3$, các số có thể thêm vào là $2$ hoặc $5$.
			\item Với mỗi số chia $3$ dư $2$ thì ta có $2$ cách  thêm số vào cuối để được $1$ số có $4$ chữ số và chia hết cho $3$, các số có thể thêm vào là $1$ hoặc $4$.
		\end{itemize}
		Do đó số các số có $4$ chữ  số và chia hết cho $3$ là $2x+2y+2z=2\cdot 180=360$.\\
		Suy  ra $n(A)=\mathrm{C}_{360}^1=360$.\\
		Vậy $\mathrm{P}(A)=\dfrac{n(A)}{n(\Omega)}=\dfrac{360}{1080}=\dfrac{1}{3} \approx 0{,}33$. 
	 }
\end{bt}

% nguồn: MaHoaNH_TraLoiNgan_Lop10
\begin{bt}%[0D0V2-7]%[ma:1002e36f]
	Chọn ngẫu nhiên 2 số trong tập hợp $X=\{1; 2; 3; \ldots; 50\}$. Tính xác suất của biến cố $B$: \lq\lq Trong hai số được chọn có một số  lớn hơn $25$, số còn lại nhỏ hơn hoặc bằng $25$\rq\rq. (kết quả làm tròn đến hàng phần trăm)
	 \shortans{0,51}
	\loigiai{ 
		Số phần tử của không gian mẫu là  $n(\Omega)=\mathrm{C}_{50}^2=1\,225$.\\
		Số cách chọn một số lớn  hơn $25$ là $25$ cách.\\
		Số cách  chọn số còn lại nhỏ hơn hoặc bằng $25$ là $25$ cách.\\
		Do  đó, ta có $n(B)=25 \cdot 25=625$.\\
		Xác suất  của biến cố $B$ là $\mathrm{P}(B)=\dfrac{n(B)}{n(\Omega)}=\dfrac{625}{1225}=\dfrac{25}{49}\approx0{,}51$.\\
	} 
\end{bt}

% nguồn: MaHoaNH_TraLoiNgan_Lop10
\begin{bt}%[0D0V2-7]%[ma:fc00f110]
	Gọi  $S$ là tập hợp tất cả các số tự nhiên gồm $4$ chữ số khác nhau được chọn từ tập hợp $A=\{1;2;3;4;5;6\}$. Chọn ngẫu nhiên một số từ tập hợp $S$. Tính xác suất để số được chọn có $2$ chữ số chẵn và $2$ chữ số lẻ.
	 \shortans{0,6}
	\loigiai{ 
		Từ  tập $A=\{1;2;3;4;5;6\}$ ta lập được $\mathrm{A}_6^4=360$ số tự nhiên có $4$ chữ số khác nhau.\\
		Phép thử ``chọn một số trong $360$ số đó'' cho ta số phần tử không gian  mẫu là $n(\Omega)=360$.\\
		Xét biến cố $X$: ``chọn được số có $2$ chữ số  chẵn và $2$ chữ số lẻ''.\\
		Số kết quả thuận lợi cho  $X$ là $n(X)=\mathrm{C}_3^2\cdot\mathrm{C}_3^2\cdot 4!=216$.\\
		Vậy  xác suất của $X$ là $\mathrm{P}(X)=\dfrac{216}{360}=\dfrac{3}{5}=0,6$.
	 }
\end{bt}

% nguồn: MaHoaNH_TraLoiNgan_Lop10
\begin{bt}%[0D0V2-7]%[ma:2430fac6]
	Có $26$ tấm thẻ được đánh số từ $1$ đến $26$. Rút ngẫu nhiên cùng một lúc $3$ tấm thẻ. Tính xác suất sao cho tích của ba số trên $3$ tấm thẻ là một  số chẵn (\textit{Kết quả làm tròn đến chữ số thập phân thứ hai sau dấu phẩy}).
	\shortans{0,89} 
	\loigiai{Số phần tử của không gian mẫu (số  cách rút $3$ thẻ bất kì trong $26$ thẻ) là $\mathrm{C}_{26}^3=2\,600$.\\
		Gọi $A$ là biến cố \lq\lq Rút được $3$  tấm thẻ mà tích các số trên chúng là một số chẵn\rq\rq.
		Ta có  $3$ trường hợp:
		\begin{itemize}
			\item \textbf{Trường hợp 1.} Rút được cả $3$ thẻ  đều là số chẵn.
			\item \textbf{Trường hợp 2.} Rút được $2$  thẻ ghi số chẵn và một thẻ ghi số lẻ.
			\item \textbf{Trường hợp 3.} Rút được $1$ thẻ ghi số  chẵn và hai thẻ ghi số lẻ.
		\end{itemize}
		 Vậy $n(A)=\mathrm{C}_{13}^3+\mathrm{C}_{13}^1 \cdot \mathrm{C}_{13}^2+\mathrm{C}_{13}^2 \cdot \mathrm{C}_{13}^1=2\,314$.\\
		 Xác xuất cần tìm là $\mathrm{P}(A)=\dfrac{n(A)}{n(\Omega)}=\dfrac{2\,314}{2\,600}=0{,}89$.
	} 
\end{bt}

% nguồn: MaHoaNH_TraLoiNgan_Lop10
\begin{bt}%[0D0V2-7]%[ma:c058626e]
	Có $20$ tấm thẻ được đánh số từ $1$ đến $20$. Chọn ngẫu nhiên ra $8$ tấm  thẻ, tính xác suất để có $3$ tấm thẻ mang số lẻ, $5$ tấm thẻ mang số chẵn trong đó chỉ có đúng $1$ tấm thẻ mang số chia hết cho $10$ (\textit{Kết quả làm tròn đến hàng phần trăm}).
	\shortans{0,13} 
	\loigiai{Không gian mẫu là cách chọn  $8$ tấm thẻ trong $20$ tấm thẻ.\\
		Suy ra số phần tử  của không mẫu là $n(\Omega)=\mathrm{C}_{20}^{8}$.\\
		Gọi $A$ là biến cố \lq\lq $3$ tấm thẻ mang số lẻ, $5$ tấm thẻ mang số chẵn trong đó chỉ  có đúng $1$ tấm thẻ mang số chia hết cho $10$\rq\rq. \\Để tìm số phần tử của $A$ ta làm như sau:
		\begin{itemize}
			\item Đầu tiên chọn $3$ tấm thẻ  trong $10$ tấm thẻ mang số lẻ, có $\mathrm{C}_{10}^3$ cách.
			\item Tiếp theo chọn $4$ tấm thẻ trong $8$ tấm thẻ mang số chẵn (không chia hết cho $10$), có $\mathrm{C}_{8}^4$  cách.
			\item Sau cùng ta chọn $1$  trong $2$ tấm thẻ mang số chia hết cho $10$, có $\mathrm{C}_{2}^1$ cách.
		\end{itemize}
		Suy ra số phần tử của biến cố $A$ là  $n(A)=\mathrm{C}_{10}^3 \cdot \mathrm{C}_{8}^4 \cdot \mathrm{C}_{2}^1$.\\
		 Vậy xác suất cần tính $\mathrm{P}(A)=\dfrac{\mathrm{C}_{10}^3 \cdot \mathrm{C}_{8}^4 \cdot \mathrm{C}_{2}^1}{\mathrm{C}_{20}^{8}}=\dfrac{560}{4\,199} \approx 0{,}13$
	} 
\end{bt}

% nguồn: MaHoaNH_TraLoiNgan_Lop10
\begin{bt}%[0D0V2-7]%[ma:2ccf1dc8]
	Gọi $S$ là tập hợp tất cả các số tự nhiên có $4$ chữ số đôi một khác nhau và  các chữ số thuộc tập hợp $\{1;2;3;4;5;6;7\}$. Chọn ngẫu nhiên một số thuộc $S$, xác suất để số đó không có hai chữ số liên tiếp nào cùng chẵn bằng bao nhiêu? (kết quả làm tròn đến hàng phần trăm)
	\shortans{0,54} 
	\loigiai{ 
		Số phần tử của tập  $S$ là $\mathrm{A}_7^4=840$ số.\\
		 Số phần tử không gian mẫu là $n(\Omega)=\mathrm{C}_{840}^1=840$ số.\\
		Gọi biến cố $A\colon $ \lq\lq không có  hai chữ số liên tiếp nào cùng chẵn\rq\rq.\\
		\textbf{Trường hợp 1}. Số  đó chỉ không có chữ số chẵn nào\\
		Trường hợp  này có $4!=24$ số.\\
		\textbf{Trường hợp 2}. Số đó chỉ có  $1$ chữ số chẵn\\
		Trường hợp này có $3\cdot\mathrm{C}_4^3\cdot4!=288$  số.\\
		\textbf{Trường hợp 3}. Số đó chứa đúng  $2$ số chẵn.\\
		Trường hợp này có 3 cách sắp xếp:(chẵn, lẻ, chẵn, lẻ) (lẻ,  chẵn, lẻ, chẵn) và (chẵn, lẻ, lẻ, chẵn) \\
		Với mỗi cách sắp xếp  thì có $\mathrm{A}_3^2\cdot\mathrm{A}_4^2=72$ số.\\
		Do đó $n(A)=24+288+3\cdot72=528$  số.\\
		 Vậy $\mathrm{P}(A)=\dfrac{n(A)}{n(\Omega)}=\dfrac{528}{840}=\dfrac{22}{35}$.
	 }
\end{bt}

% nguồn: MaHoaNH_TraLoiNgan_Lop10
\begin{bt}%[0D0C2-7]%[ma:ee551dcc]
	Cho tập hợp $X=\{0;1;2;3;4;5;6;7;8\}$. Gọi $S$ là tập hợp tất cả các số tự nhiên có $5$ chữ số đôi một khác nhau được lập từ tập hợp $X$. Lấy ngẫu nhiên một số từ $S$. Biết xác suất để số lấy được  chia hết cho $5$ bằng $\dfrac{a}{b}$ với $\dfrac{a}{b}$ là phân số tối giản; $a$, $b$ nguyên dương. Tính giá trị của biểu thức $a+b$.
	\shortans{79} 
	 \loigiai{
		Gọi số cần lập là $\overline{abcde}$ với $a,b,c,d,e\in  X$ và khác nhau đôi một.\\
		Chọn chữ  số $a$ có $8$ cách chọn vì $a\ne 0$.\\
		Chọn $4$ chữ số khác chữ số $a$ đã chọn từ tập  hợp $X$ và sắp xếp vào $4$ vị trí của $b$, $c$, $d$, $e$ có $\mathrm{A}_8^4=1680$ cách.\\
		Suy ra, có  $8\cdot 1680=13440$ số.\\
		Do  đó, $n(\Omega)=13440$.\\
		 Gọi $A$ là biến cố \lq\lq số được chọn chia hết cho $5$\rq\rq.
		\begin{itemize}
			\item $\overline{abcd0}$ là số chia hết cho  $5$.\\
			Chọn $4$ chữ số  khác từ tập hợp $X$ và sắp xếp vào $4$ vị trí của $a$, $b$, $c$, $d$ có $\mathrm{A}_8^4=1680$ cách
			\item $\overline{abcd5}$ là số chia  hết cho $5$.\\
			 Chọn chữ số $a$ có $7$ cách chọn vì $a\ne 0$ và $a\ne 5$.\\
			Chọn $3$ chữ số khác chũ số $a$ đã chọn và chữ số $5$ từ tập hợp $X$ và sắp xếp vào $3$ vị trí của $b$,  $c$, $d$ có $\mathrm{A}_7^3=210$ cách.\\
			Suy ra có  $7\cdot 210=1470$
		\end{itemize}
		Do đó,  $n(A)=3150$.\\
		 Vậy xác suất cần tìm là $\mathrm{P}(A)=\dfrac{n(A)}{n(\Omega)}=\dfrac{3150}{13440}=\dfrac{15}{64}$.
	 }
\end{bt}

% nguồn: MaHoaNH_TraLoiNgan_Lop10
\begin{bt}%[0D0C2-7]%[ma:2bea897c]
	Gọi $X$ là tập hợp các số tự nhiên có $3$ chữ số phân biệt được lập từ các số $0$, $1$, $2$,  $3$, $4$, $5$. Chọn ngẫu nhiên một số từ $X$. Tính xác suất để có thể chọn được một số không chia hết cho $3$.
	\shortans{0,6} 
	\loigiai{ 
		Số lượng  phần tử của tập hợp $X$ là $5\cdot5\cdot4=100$.\\
		 Chọn $1$ số trong tập $X$ thì số phần tử của không gian mẫu là $n(\Omega)=100$.\\
		Gọi $A$ là biến cố \lq\lq Chọn được một  số không chia hết cho $3$\rq\rq\,.\\
		Suy  ra $\overline{A}$ là biến cố \lq\lq Chọn được một số chia hết cho $3$\rq\rq\,.\\
		Dấu hiệu chia hết cho $3$ là tổng các chữ số là một số chia hết cho $3$ nên ta xét  các trường hợp sau:
		\begin{itemize}
			\item Tổng các chữ số bằng $12\Rightarrow$ xét tập  các chữ số $\{3;4;5\}$ lập đươc $3!=6$ số chia hết cho $3$.
			\item Tổng  các chữ số bằng $9$:
			\begin{itemize}
				\item Xét tập các chữ số $\{0;4;5\}$ lập  đươc $2\cdot2!=4$ số chia hết cho $3$.
				\item Xét tập các chữ số  $\{1;3;5\}$ lập đươc có $3!=6$ số chia hết cho $3$.
				\item Xét tập các  chữ số $\{2;3;4\}$ lập đươc có $3!=6$ số chia hết cho $3$.
			\end{itemize}
			 \item Tổng các chữ số bằng $6$:
			\begin{itemize}
				\item Xét tập các chữ  số $\{0;1;5\}$ lập đươc có $2\cdot2!=4$ số chia hết cho $3$.
				 \item Xét tập các chữ số $\{0;2;4\}$ lập đươc có $2\cdot2!=4$ số chia hết cho $3$.
				 \item Xét tập các chữ số $\{1;2;3\}$ lập đươc có $3!=6$ số chia hết cho $3$.
			\end{itemize}
			\item Tổng các chữ số bằng $3\Rightarrow$ xét tập các  chữ số $\{0;1;2\}$: có $2\cdot2!=4$ số chia hết cho $3$.
		\end{itemize}
		Theo quy tắc cộng có  ta có $n\left(\overline{A}\right)=6+4+6+6+4+4+6+4=40$.
		\[ 
		\mathrm{P}\left(\overline{A}\right)=\dfrac{40}{100}=0,4\Rightarrow\mathrm{P}(A)=1-0{,}4=0{,}6. 
		 \]
		Vậy xác suất để có thể chọn được  một số không chia hết cho $3$ là $0{,}6$.
	} 
\end{bt}


%%%%%%%%%%%%%%%% TỰ LUẬN %%%%%%%%%%%%%%%%

% nguồn: MaHoaNH_TuLuan_Lop10
\begin{bt}%[0D0H2-7]%[ma:26ee544e]
	Có ba chiếc bình khác  nhau, mỗi bình chứa $3$ tấm thẻ được đánh số $1$, $2$, $3$. Từ mỗi bình rút ngẫu nhiên một tấm thẻ. Mô tả không gian mẫu và tính xác suất của các biến cố
	\begin{enumerate}
		\item $A$:\lq \lq Tổng  số thẻ ghi trên các tấm thẻ bằng $6$ \rq \rq.
		\item $B$: \lq \lq Tích  $3$ số ghi trên $3$ tấm thẻ là số lẻ\rq \rq.
	\end{enumerate}
	\loigiai{Không gian mẫu $\Omega = \left\{ \left(i,j,k\right), \; 1  \le i,j,k \le 3, \; i,j,k \in \mathbb{N}\right\}$. Do đó $n(\Omega)=27$.
		\begin{enumerate}
			\item  Ta có $A=\left\{(1,2,3);(1,3,2);(2,1,3),(2,3,1);(3,1,2);(3,2,1);(2,2,2)\right\}$.\\
			Xác suất của biến  cố $A$ là $\mathrm{P}(A)=\dfrac{n(A)}{n(\Omega)}=\dfrac{7}{27}$.
			\item Ta có  $B=\left\{(1,1,1);(1,1,3);(1,3,1),(1,3,3);(3,1,1);(3,1,3);(3,3,1);(3,3,3)\right\}$.\\
			Xác  suất của biến cố $B$ là $\mathrm{P}(B)=\dfrac{n(B)}{n(\Omega)}=\dfrac{8}{27}$.
	\end{enumerate}}
\end{bt}

% nguồn: Ok_TuanLuan_Toan11_Lan1
\begin{bt}%[0D0H2-7]%[ma:cac6df72]
	Một hộp có $20$ chiếc  thẻ cùng loại, mỗi thẻ được ghi một trong các số $1; 2; 3; \ldots; 20$, hai thẻ khác nhau thì ghi hai số khác nhau. Rút ngẫu nhiên $1$ chiếc thẻ trong hộp. Xác suất để rút được thẻ ghi số chia hết cho $3$ hoặc chia hết cho $7$ bằng bao nhiêu?
	\loigiai{ 
		Gọi $A$ là biến cố \lq\lq Rút $1$ thẻ và số ghi trên thẻ chia hết cho $3$ hoặc chia hết cho $7$\rq\rq.\\ 
		Rút  một thẻ từ $20$ thẻ nên số phần tử không gian mẫu là $n(\Omega)=20$ (cách).\\
		\textbf{Trường hợp 1.} Rút được $1$ tấm thẻ  và số ghi trên thẻ chia hết cho $3$ thì số trên thẻ thuộc $\{3;6;9;12;15;18\}$ có $6$ (cách).\\
		\textbf{Trường hợp 2.} Rút được $1$ tấm thẻ và số ghi trên thẻ chia hết  cho $7$ thì số trên thẻ thuộc $\{7;14\}$ có $2$ cách.
		Vậy xác  suất của biến cố $A$ là $\mathrm{P}(A)=\dfrac{6+2}{20}=\dfrac{3}{5}$.
	} 
\end{bt}

% nguồn: Ok_TuanLuan_Toan11_Lan2
\begin{ex}%[0D0H2-7]%[ma:bd6ee85a]
	Một hộp có $20$ chiếc thẻ cùng loại, mỗi thẻ được  ghi một trong các số $1; 2; 3; \ldots; 20$, hai thẻ khác nhau thì ghi hai số khác nhau. Rút ngẫu nhiên $1$ chiếc thẻ trong hộp. Xác suất để rút được thẻ ghi số chia hết cho $3$ hoặc chia hết cho $7$ bằng bao nhiêu?
	 \loigiai{
		Gọi $A$ là biến cố \lq\lq Rút $1$ thẻ và số ghi trên thẻ  chia hết cho $3$ hoặc chia hết cho $7$\rq\rq.\\
		Rút  một thẻ từ $20$ thẻ nên số phần tử không gian mẫu là $n(\Omega)=20$ (cách).\\
		\textbf{Trường hợp 1.} Rút được $1$ tấm thẻ và số ghi trên thẻ chia hết cho $3$ thì số  trên thẻ thuộc $\{3;6;9;12;15;18\}$ có $6$ (cách).\\
		\textbf{Trường hợp 2.} Rút được $1$  tấm thẻ và số ghi trên thẻ chia hết cho $7$ thì số trên thẻ thuộc $\{7;14\}$ có $2$ cách.
		Vậy xác suất của biến cố $A$  là $\mathrm{P}(A)=\dfrac{6+2}{20}=\dfrac{3}{5}$.
	} 
\end{ex}
